\SourceLeaf{436}%
\Needspace{7\baselineskip}\begin{center}13.\end{center}
\noindent\hspace*{1.2em}Einem jeden Punkte einer krummen Linie auf der krummen Fläche entsprechen auf der Kugelfläche nach unserer Ansicht bereits zwei Punkte, nemlich der Punkt $\lambda$, der die Richtung des Linienelements und der Punkt $L$, der die Richtung der Normale auf die Fläche vorstellt. Beide sind offenbar $90^{0}$ von einander entfernt. Bei unserer frühern Untersuchung (Art. $9$), wo [wir] die krumme Linie in Einer Ebene liegend vorausgesetzt hatten, hatten wir noch \textls[100]{zwei} andere Punkte auf der Kugelfläche, nemlich $\mathfrak{L}$, der die Richtung der Normale auf die Ebene, und $\lambda'$, der die Richtung der Normale auf das Element der Curve in der Ebene vorstellte; hierbei war also $\mathfrak{L}$ ein fester Punkt und $\lambda$, $\lambda'$ immer in Einem grössten Kreise, dessen Pol $\mathfrak{L}$. Indem wir jetzt die Betrachtung generalisiren, wollen wir zwar die Bezeichnungen $\mathfrak{L}$, $\lambda'$ beibehalten, müssen aber ihre Bedeutung aus einem allgemeinern Gesichtspunkt auffassen. Indem nemlich die Curve $s$ beschrieben wird, werden die Punkte $L,\lambda$ auf der Hülfskugelfläche auch krumme Linien beschreiben, die, allgemein zu reden, keine grössten Kreise mehr sind. Mit dem Element der zweiten Linie parallel ziehe man einen Radius der Hülfskugel zum Punkt $\lambda'$, statt dessen wir aber den gegenüberliegenden Punkt wählen, wenn jener mehr als $90^{0}$ von $L$ entfernt ist; im ersten Fall betrachten wir das Element an $\lambda$ als positiv, im andern als negativ; endlich sei $\mathfrak{L}$ der Punkt der Hülfskugel, der von $\lambda$ und $\lambda'$ zugleich $90^{0}$ entfernt ist und zwar so, dass $\lambda,\lambda',\mathfrak{L}$ in derselben Ordnung liegen wie $(1),(2),(3)$.

Die Coordinaten der vier Punkte auf der Hülfskugel relativ gegen den Mittelpunkt sind für
\[
\begin{array}{cccc}
L&X&Y&Z\\
\lambda&\xi&\eta&\zeta\\
\lambda'&\xi'&\eta'&\zeta'\\
\mathfrak{L}&\alpha&\beta&\gamma.
\end{array}
\]
Jeder dieser $4$ Punkte beschreibt also auf der Hülfskugel selbst eine Linie, deren Elemente wir durch $dL,d\lambda,d\lambda',d\mathfrak{L}$ ausdrücken wollen. Wir haben also
